\section{A redundant truth equation: an exact 667-bit equational packet}
\label{sec:e667}

The inspected incoming packet \(\mathcal E_{692}\) has an equation that can be
deleted without changing its equational theory.  The new packet
\(\mathcal E_{667}\) has seven equations, nine operation symbols, four variable
slots, and 150 equation-term nodes.  Its exact serialization is 667 logical
bits.  This result concerns a finite equational consistency problem; it does
not give a Turing machine with 667, 237, or any other specified number of states.

The exact input is the \texttt{E692.json} member of
\path{docs/incoming/ZFC_Sub237_Research_2026-10-19.zip} at ProveIt commit
\path{1b27f052d770f4c7d435e82e779bf21f5a1f66ed}.  Its raw SHA-256 is
\path{0657cefef8f41bd7e463eb0adcd5a44b118de62ec852729202d3ecc4b6942d43}.
The original source's descriptor, complemented binder, structural witnesses,
selector convention, and bit codec are retained.  Only the equation
\begin{equation}
                         d=N(d,N(d,d))
\label{eq:e667-deleted}
\end{equation}
is deleted.

\paragraph{Earlier truth elimination and the changed premise.}
Source 104 already deleted this explicit truth equation from the older
E946 packet, obtaining E921.  Its proof used the separate arbitrary-witness
equation of the ordinary existential interface.  E692 has a complemented
binder instead, so that literal proof premise is absent here.  The present
contribution is the new closure-based deduction below and its application
to the combined E692 packet, not the general idea of deleting a redundant
truth equation.  The prior argument is retained in ProveIt at
\texttt{data/104-res238-formalisms-truth\_redundancy.md} within the
\texttt{zfc-closure-kernels} reports.

\subsection{The precise Boolean premise}

Both packets retain the single axiom
\begin{equation}
 N\bigl(N(x,N(N(y,x),x)),N(y,N(z,x))\bigr)=y.
\label{eq:e667-sh1}
\end{equation}
This is exactly Sh1 in Theorem 2 of McCune, Veroff, Fitelson, Harris, Feist,
and Wos, \emph{Short Single Axioms for Boolean Algebra}, Journal of Automated
Reasoning 29 (2002), 1--16, DOI
\href{https://doi.org/10.1023/A:1020542009983}{10.1023/A:1020542009983}.
The authors give a finite 66-step derivation of the Sheffer basis from Sh1;
the \href{https://www.cs.unm.edu/~mccune/papers/basax/v12.pdf}{author-hosted
paper} and \href{https://www.cs.unm.edu/~mccune/papers/basax/Sh-1.proof}{plain
proof} provide the primary source.
Consequently the reduct of every Sh1 algebra is a Boolean algebra in which
\(N(a,b)=\neg(a\wedge b)\), after choosing the corresponding Boolean dual if
necessary.  No assumption concerning the distinguished constant \(d\) is
needed for this interpretation.  In particular, the deleted equation is
not used to obtain Boolean algebra.

\subsection{A short equational lemma}

For an arbitrary Boolean-algebra element \(u\), define the term
\[
                   t(u)=N(u,N(u,u)).
\]
It denotes the Boolean top, independently of \(u\).  For arbitrary
\(d,k,h,u\), put
\[
              F=N\bigl(N(t(u),N(d,k)),N(d,h)\bigr).
\]
Then the following is an identity of Boolean algebras:
\begin{equation}
                         N(F,N(d,d))=d.
\label{eq:e667-identity}
\end{equation}
Indeed, expansion using NAND gives
\[
 F=\neg d\vee\neg k\vee h,
 \qquad F\wedge\neg d=\neg d,
 \qquad N(F,N(d,d))=\neg\neg d=d.
\]
These are finite equational manipulations in the Boolean reduct.  They do
not use an ultrafilter, an infinite meet, a two-element model assumption,
or the equation \(d=1\).  Composing with the cited Sh1 derivation makes
\eqref{eq:e667-identity} an equational consequence of the single retained
axiom \eqref{eq:e667-sh1}.

It follows immediately that an equation \(d=F\) forces
\[
 N(d,N(d,d))=N(F,N(d,d))=d,
\]
which is exactly the omitted equation, with its sides exchanged.

\subsection{The retained closure has exactly the required form}

Write \(\varepsilon\) for the packet's binary operation
\(p_\varepsilon\), and retain \(p_E\) as its other binary predicate-valued
operation.  All notation in the next display abbreviates literal terms;
no new operation is added to the packet:
\begin{align*}
 r(b,z)&=\varepsilon\bigl(P(P(b,z),P(Q(b),b)),V(b,d)\bigr),\\
 k(b,z)&=N\bigl(p_E(A(b,z),P(b,z)),N(d,p_E(A(b,z),z))\bigr),\\
 c(b,x,y)&=\varepsilon\bigl(P(y,P(x,b)),V(b,d)\bigr),\\
 u(b,y)&=p_E(y,Q(b)),\\
 v(b,y)&=\varepsilon\bigl(P(y,P(P(b,y),b)),V(b,d)\bigr),\\
 a(x)&=p_E(W(x,d),x),\qquad b_0(b,x)=p_E(x,Q(b)),\\
 h(b,x,y)&=N\Bigl(
   N\bigl(N(d,u(b,y)),N(c(b,x,y),N(a(x),N(d,b_0(b,x))))\bigr),
   N(u(b,y),v(b,y))\Bigr).
\end{align*}
The last retained equation of \(\mathcal E_{692}\) is literally
\begin{equation}
 d=N\Bigl(N\bigl(r(b,z),N(d,k(b,z))\bigr),N(d,h(b,x,y))\Bigr).
\label{eq:e667-closure}
\end{equation}
There is no use of the intended set-theoretic semantics in identifying
this syntax.  In particular, \(h\) can be treated as a completely arbitrary
Boolean-algebra element in the lemma.

The other two retained equations needed for the new argument are
\begin{align}
                   V(P(x,s),d)&=x,\label{eq:e667-head}\\
 \varepsilon(s,P(a,b))&=N(\varepsilon(s,a),\varepsilon(s,b)).
 \label{eq:e667-class}
\end{align}
Let the \emph{ground class term} and descriptor be
\[
             T=P(d,P(d,d)),\qquad B=P(T,d).
\]
For every environment \(S\), two instances of
\eqref{eq:e667-class} give
\begin{equation}
 \varepsilon(S,T)
   =N\bigl(\varepsilon(S,d),N(\varepsilon(S,d),\varepsilon(S,d))\bigr)
   =t(\varepsilon(S,d)).
\label{eq:e667-tautology}
\end{equation}
This works for every meaning of the class \(d\); it is not an assumption
that \(d\) denotes a true class.  By \eqref{eq:e667-head}, \(V(B,d)=T\),
so, on writing \(S_z=P(P(B,z),P(Q(B),B))\),
\[
                     r(B,z)=t(\varepsilon(S_z,d)).
\]
Substitute \(b=B\) into \eqref{eq:e667-closure}.  Its right-hand side is
now exactly the lemma's \(F\), with
\[
 u=\varepsilon(S_z,d),\quad k=k(B,z),\quad h=h(B,x,y).
\]
Equation \eqref{eq:e667-identity} and congruence give
\eqref{eq:e667-deleted}.  If a closed derivation is desired, set
\(x=y=z=d\) as well: one ground instance of the closure row, the head
row, two class-NAND instances, and the Boolean identity suffice.

\subsection{Exact equivalence and its scope}

Every equation of \(\mathcal E_{667}\) is an equation of
\(\mathcal E_{692}\).  The preceding finite argument derives the one
omitted equation in the other direction.  Therefore, in the same
signature, the two presentations have exactly the same equational
consequences and exactly the same models, including the one-element model.
In particular,
\begin{equation}
 \mathcal E_{667}\vdash N(d,d)=d
 \quad\Longleftrightarrow\quad
 \mathcal E_{692}\vdash N(d,d)=d.
\label{eq:e667-equivalence}
\end{equation}
This is stronger than merely having the same consistency strength.
Its proof does not depend on the raw-EPSC calibration.

The incoming E692 report supplies the separate ordinary relative model
argument connecting that packet to E713 and raw EPSC.  Composing
\eqref{eq:e667-equivalence} with that explicitly inherited calibration gives
the intended ZFC-metatheoretic equivalence between the nonderivability of
the displayed ground query and \(\operatorname{Con}(\mathrm{ZFC})\).
This work does not replace the inherited model argument or claim a new
Lean/Rocq or internally serialized PA proof.  It also supplies no literal
binary machine for the new equational presentation.

\subsection{Reproduction and accounting}

\texttt{build\_e667.py} checks the exact input hashes and reconstructs the
692-bit input using an independently written encoder.  It identifies the
literal retained premises above and the exact shape of
\eqref{eq:e667-closure}, removes just equation index 1, and writes the new
packet, expanded equations, syntax tree, and verification receipt.
An independently written strict decoder checks the generated syntax and
endpoints.  The new Boolean identity is normalized as a squarefree
polynomial over \(\mathbb F_2\), using \(N(a,b)=1+ab\); its two sides
normalize to the same polynomial \(d\).  This symbolic check is valid in
every Boolean ring, with no finite cardinality bound.  The use of the
published Sh1 theorem remains explicit.

The deleted equation has six term nodes and one equation marker.  The
common opcode width remains four bits, hence the saving is
\[
                         6\cdot4+1=25\text{ bits}.
\]
The exact total is
\[
 \underbrace{43}_{\text{signature and variable header}}
 +\underbrace{(7+4\cdot150+1)}_{\text{seven equations and terminator}}
 +\underbrace{16}_{\text{two ground endpoints}}
 =667.
\]
The physical transport file has 84 bytes with five trailing zero padding
bits.  Its logical length is recorded separately; the padding is not part
of the packet.  Reproduce with
\begin{verbatim}
python3 -B build_e667.py --output /tmp/e667-replay
\end{verbatim}

The binary-state objective remains separate: the new theorem saves one
axiom and six emitter-term nodes in an alternative frontend, but a complete
equational enumerator, writer, initializer, and literal table must still be
accounted before any state reduction can be asserted.
